\begin{abstract}
Monitoring forest health at scale requires processing decades of satellite imagery across millions of pixels. This literature survey examines the intersection of change-detection algorithms and scalable geospatial data systems that together enable large-area forest disturbance analysis. We review statistical time-series methods such as BFAST and Adaptive EWMA, spectral and species-specific approaches using Sentinel-2 and Landsat, unsupervised deep-learning models, and the distributed processing frameworks including RDPro, Apache Sedona, and space-time data cubes that make continent-scale computation feasible. We also consider the accuracy of widely used global forest products like GLAD in temperate climates. The survey highlights that scalable analysis depends not only on algorithmic innovation but equally on efficient storage layouts, spatiotemporal indexing, and cloud-native data formats.

\end{abstract}